\scanpage{48}
\mspage{46}

rie finie. (Je ne me rends pas bien compte si \unsure{sa|ça} arrive très souvent\ldots\ \deleted{\unsure{)}} Mais on ne peut fonder une théorie générale des dérivateurs et des catégories de modèles ``dérivables'', sur une hypothèse de ce genre.)

Mais supposons que le foncteur $H$ envisagé se factorise même par \unsure{Hot}, i.e. qu'il transforme quasi-iso. en bijections. On a vu que cela équivaut à la condition qu'il puisse se représenter par un système inductif sur $\mathbb N$ formé de catégories de Thomason-Kan.

Il doit être facile de plus, si on y tient, de s'arranger pour que les foncteurs de transition soient des mono et pl. fid., pour avoir une $\underline{\lim}$ \unsure{sympa}. Ce sera bien \unsure{sur|sûr} une catégorie de Thomason-Kan.

Cette catégorie est-elle déterminée par le foncteur $H:\mathrm{Cat\,hot}^{\unsure{o|0}}\unsure{\relbar|\to}\mathrm{Ens}$, à iso unique près \unsure{!|à} dans Cat hot ? Bien sûr, il n'est pas question qu'elle représente le foncteur $H$, \uline{sauf} si \inserted{celui-ci est bel et bien représentable, i.e.} le \deleted{système inductif} \inserted{ind objet} $\underline{\underline K}$ de Cat hot est essentiellement constant (ce qui ne doit pratiquement jamais être le cas). Notons déjà

% archive sheet number: 47
